\subsection{Finite arithmetic matrices over an arbitrary number field}

\begin{statement}{C.14}{Proved}{An explicit finite Galois construction of the retained arithmetic groups}
Fix a number field $L\subset\overline{\mathbb Q}$ and put
$S=\{v:v\mid3\}\cup\{v:v\mid\infty\}$.
For $n\ge1$ write
\[
 A_n=\mathcal K/3^n,\qquad
 A_n^D=\operatorname{Hom}(A_n,\mu_{3^n}),\qquad
 F_n=L(J_{E,D}[3^n],\mu_{3^n}),\qquad
 d_n=3^{44n}[F_n:L].                                      \tag{C.14.1}
\]
The homomorphisms in the dual have their usual contragredient
Galois action. Let $\Omega_n$ be the compositum, inside the
fixed algebraic closure, of all finite Galois extensions $E/L$
containing $F_n$, unramified outside $S$, and satisfying
$[E:L]\le d_n$. Put $G_n=\operatorname{Gal}(\Omega_n/L)$.
At each $v\mid3$ choose a prolongation to $\Omega_n$ and
write $D_{n,v}\le G_n$ for its decomposition subgroup.
These are finite fields and finite groups, obtained by the
bounded enumeration in the proof.

For a finite $G_n$-module $A$, use the finite matrices
\[
 \begin{split}
  (\delta^0 b)(g)&=g b-b,\\
  (\delta^1 z)(g,h)&=g z(h)-z(gh)+z(g).
 \end{split}                                               \tag{C.14.2}
\]
Define the following finite abelian groups by kernels and
images of these matrices:
\[
 U_n=\ker\bigl(\delta^1:\operatorname{Map}(G_n,A_n)
                   \to\operatorname{Map}(G_n^2,A_n)\bigr)
                      /\delta^0 A_n,                       \tag{C.14.3}
\]
and
\[
\begin{split}
 C_n={}&\{(z,(b_v)_{v\mid3}):
       \delta^1z=0,\ b_v\in A_n^D,\\
     &\hspace{26mm}z(g)=g b_v-b_v
             \text{ for }g\in D_{n,v}\}\big/
        \{(\delta^0 b,(b)_{v\mid3}):b\in A_n^D\}.
\end{split}                                                 \tag{C.14.4}
\]
In (C.14.4), $z$ ranges over all functions $G_n\to A_n^D$.
There are canonical identifications
\[
\begin{split}
 H^1(G_{L,S},A_n)&\simeq U_n,\qquad
 H^2(G_{L,S},A_n)\simeq C_n^\vee,\\
 H^1_{\rm cts}(G_{L,S},\mathcal K)&\simeq\varprojlim_n U_n,\qquad
 H^2_{\rm cts}(G_{L,S},\mathcal K)
       \simeq\varprojlim_n C_n^\vee .
\end{split}                                                 \tag{C.14.5}
\]
Here $C_n^\vee=\operatorname{Hom}(C_n,\mathbb Q/\mathbb Z)$;
the transition maps are induced by coefficient reduction
and its Cartier dual. Thus the arithmetic terms of C.12
are specified by finite arithmetic fields and finite matrices
at every precision. The construction applies to every $L$;
it does not assume that either displayed inverse limit is zero.
\end{statement}
\begin{proof}
The reduction $A_n$ is a submodule of
$\bigwedge^2(V/3^n)$: the quotient $\bigwedge^2T$ in the
definition of $\mathcal K$ is torsion-free. The finite group
scheme $J_{E,D}[3^n]$ therefore splits its action. Adjoining
$\mu_{3^n}$ also splits the Cartier dual. Both coefficients
are unramified outside $S$, because the original generalized
Jacobian and its three-adic representation are unramified
outside three. All the fields in (C.14.1) are algebraic data
of the specified curve. For example, multiplication by $3^n$
on its framed divisor classes constructs $J_{E,D}[3^n]$;
Riemann--Roch and factorization give its finite defining
equations and their splitting field. No presentation of an
unknown infinite Galois group is an input to the construction.

We give a bound for the field enumeration. For any integer
$d\ge1$, every extension $E/L$ of degree at most $d$
unramified outside $S$ has
\[
 |\operatorname{disc}(E)|\le
 B_L(d):=|\operatorname{disc}(L)|^d
       \prod_{v\mid3}(Nv)^{d(1+e_v\lfloor\log_3 d\rfloor)},
                                                               \tag{C.14.6}
\]
where $e_v$ is the ramification index of $L_v/\mathbb Q_3$.
Indeed, pass first to the maximal unramified subextension
of a completion $E_w/L_v$. Its remaining totally ramified
part has ramification index $e$ and an Eisenstein polynomial
$P$ for a uniformizer $\pi$. The terms of $P'(\pi)$ have
pairwise distinct valuations modulo $e$. Its valuation is
therefore their minimum and is at most
$e-1+e e_v v_3(e)$. Multiplying by the residue degree and
summing over $w\mid v$ bounds the relative discriminant
exponent by $d(1+e_v\lfloor\log_3d\rfloor)$. The tower
formula for the discriminant proves (C.14.6).

Hermite's proof, in the explicit form of Milne,
Theorems 8.42--8.43 \cite{CMilneANT}, bounds the conjugates
of a primitive integral element in terms of the absolute
discriminant. Consequently one can list a finite collection
of monic integral polynomials of degree at most $d[L:\mathbb Q]$
that contains a defining polynomial for every $E$ in question.
For a deliberately generous choice, a bound
$2(B_L(d)+1)$ for every conjugate suffices in that proof;
the coefficient of degree $r-i$ is then bounded by
$\binom ri[2(B_L(d)+1)]^i$. Enumerate these polynomials,
factor them, calculate rings of integers and discriminants,
and retain precisely their embeddings over $L$ that are
Galois, contain $F_n$, have the required degree, and are
unramified away from $S$. All these tests are finite exact
number-field operations. Taking their compositum constructs
$\Omega_n$. The bound is used to prove termination; no
feasibility estimate is asserted.

Let $A$ be either $A_n$ or $A_n^D$, and let
$\rho:G_{L,S}\to\operatorname{Gal}(F_n/L)$ be the quotient.
Every continuous crossed homomorphism $z:G_{L,S}\to A$
gives the continuous homomorphism
\[
       g\longmapsto(z(g),\rho(g)):
       G_{L,S}\longrightarrow A\rtimes\operatorname{Gal}(F_n/L).
                                                               \tag{C.14.7}
\]
Its kernel cuts out a Galois extension containing $F_n$
of degree at most $|A|[F_n:L]=d_n$, unramified outside $S$.
It is one of the fields used in $\Omega_n$. Thus every
crossed homomorphism factors through $G_n$.
Conversely every crossed homomorphism on $G_n$ inflates to
one on $G_{L,S}$, and the coboundaries on the two groups
are exactly the same functions $g\mapsto gb-b$.
This proves (C.14.3) and its analogue for $A_n^D$.

Use the compact-support complex
\[
 R\Gamma_c(\mathcal O_L[1/3],A_n^D)
 =\operatorname{Cone}\!\left(
 R\Gamma(G_{L,S},A_n^D)
 \longrightarrow\bigoplus_{v\in S}R\Gamma(L_v,A_n^D)
                          \right)[-1].                       \tag{C.14.8}
\]
The local complexes at infinite places here are the modified
Tate complexes. They are acyclic for these odd-primary
coefficients, including at a real place. Degree-one cycles
in (C.14.8) are exactly a global one-cocycle together with
the indicated local degree-zero trivializations. A change
by a global zero-cochain is exactly the denominator of
(C.14.4). Since the global cocycle factors through $G_n$,
each local equation holds on $G_{L_v}$ exactly when it
holds on its image $D_{n,v}$. We have proved the explicit
identification of $C_n$ with that compact-support group.

Milne's Corollary II.3.3(a),(b) \cite{MilneADT}, applied
to the locally constant sheaf $A_n^D$, identifies its
derived Cartier dual with $A_n$ and gives the perfect
pairing
\[
 H^2(G_{L,S},A_n)\times C_n
               \longrightarrow\mathbb Q/\mathbb Z.          \tag{C.14.9}
\]
The Galois-to-etale comparison is the same II.2.9 used
in C.11. This proves the second finite identification.
Every finite group in degrees one and two is finite, so
each coefficient tower is Mittag--Leffler. The continuous
cohomology inverse-limit sequences therefore have zero
first derived-limit term at finite $S$ and give the last
two identifications in (C.14.5). Coefficient naturality
of (C.14.7)--(C.14.9) specifies all transition maps.

For a small arithmetic verification of the matrix convention,
the program \texttt{c\_arithmetic\_matrices.py}
in \texttt{code/} constructs the actual
group of $\mathbb Q(\zeta_9,\sqrt[3]3)/\mathbb Q$, of order
$18$. Over $K_0$ the two Kummer classes $\zeta,3$ are
independent and span the $S$-unit classes: $3=-\zeta^2\lambda^2$.
Thus this quotient captures all one-cocycles for
$\mathbb F_3$ and $\mu_3$ over $\mathbb Q$, including
their coboundaries. Its decomposition group at three is
the whole group. Indeed a Kummer line involving $3$
has valuation nonzero modulo three over $K_{0,3}$;
the remaining line $\zeta$ gives the ramified cyclotomic
cubic extension. Hence no nontrivial cubic subextension
is unramified there. In coordinates the group law is
$(a,b,c)(a',b',c')=(a+a',b+(-1)^c b',c+c')$,
with the first two coordinates modulo three and the
last modulo two. The program verifies this law and
both coefficient actions. The respective cocycle
matrix ranks are $17$ and $16$, the coboundary ranks
are $0$ and $1$, and both first-cohomology dimensions
are one. Both compact-support dimensions are zero.
A separately labelled duplicated-subgroup fixture
checks the local-trivialization quotient; it is not
assigned to a number field. These are checks of the
finite matrices, with their complete output retained;
the proof above establishes their general arithmetic use.
\end{proof}

\begin{statement}{C.15}{Proved}{The complete number-field class-two test with its finite arithmetic evaluations}
Let $s$ be a continuous section of $N^{\rm cc}$ over $L$.
Use the actual obstruction two-cocycle $o$ and its
origin-corrected residue family $r=(r_v)_{v\nmid3,\infty}$
from A.8, restricted to $G_L$. Write
$D_v=\operatorname{Frob}_v-1$ on $\mathcal K(-1)$.
First require
\[
 y_v=-D_v^{-1}r_v\in\mathcal K(-1),\qquad
 (y_v)_v\in\mathcal R_L(\mathcal K(-1)).                    \tag{C.15.1}
\]
The actual localization homotopy-fibre cycle $(o,y)$
then supplies, at each $n$, a two-cocycle
$\xi_n\in Z^2(G_{L,S},A_n)$, by the comparison below.
For every class represented by $(z,(b_v))$ in the
explicit finite group $C_n$ of C.14, choose a continuous
two-cochain $a$ with values in $\mu_{3^n}$ satisfying
\[
                        da=\xi_n\cup z.                    \tag{C.15.2}
\]
Such a cochain always exists. Define
\[
 \ell_n(s)(z,(b_v))
   =\sum_{v\mid3}\operatorname{inv}_v
       \bigl((\operatorname{res}_v\xi_n)\cup b_v
                          -\operatorname{res}_v a\bigr).
                                                               \tag{C.15.3}
\]
The expression in each local invariant is a two-cocycle.
Formula (C.15.3) is independent of all the auxiliary
cochain choices and is an element of $C_n^\vee$.
There is an exact criterion
\[
 s\text{ lifts to }N/\gamma_3N
 \quad\Longleftrightarrow\quad
 \text{(C.15.1) holds and }
        \ell_n(s)=0\text{ for every }n\ge1.                 \tag{C.15.4}
\]
Every nonempty fibre is a torsor under
$\varprojlim_n U_n$, with $U_n$ given by the finite
matrices (C.14.3), and no further geometric-stabilizer
quotient. Together with C.13 for the initial lift,
C.14--C.15 give the number-field version requested
in Task C.1: both the residue convergence and the
arithmetic test are retained.
\end{statement}
\begin{proof}
We specify the comparison that enters (C.15.2).
Put $\mathcal A_v=\mathcal K(-1)$ and
$R^\bullet=[\mathcal R_L(\mathcal A)
                     \xrightarrow{D}\mathcal R_L(\mathcal A)]$
in degrees zero and one. The localization triangle of F.1
provides a continuous cochain model and a quasi-isomorphism
\[
 R\Gamma(G_{L,S},\mathcal K)
 \longrightarrow
 \operatorname{Cone}\!\left(
   R\Gamma(G_L,\mathcal K)\longrightarrow R^\bullet[-1]
                         \right)[-1].                       \tag{C.15.5}
\]
Choose the cone sign so that a degree-two cycle is
$(o,y)$ with $do=0$ and $r+Dy=0$. This is the sign in
(C.15.1). All restrictions, residue maps, and homotopies
are the actual maps in the localization construction;
in particular $y$ is retained as part of the cycle.
Reduction modulo $3^n$ in (C.15.5) gives its finite
coefficient comparison. A representative $\xi_n$ is
any finite-coefficient two-cocycle on the left whose
image is cohomologous to the reduction of $(o,y)$ on
the right. It exists because (C.15.5) is a
quasi-isomorphism. Its cohomology class is uniquely
determined. This prescription keeps the arithmetic
class in the homotopy fibre, including when the
finite-coefficient map from that class to global
cohomology has a kernel.

There is a finite-data method for obtaining such
representatives at any specified $n$. Continuous
finite-valued cochains factor through finite Galois
quotients. Enumerate finite $S$-ramified quotients
and their normalized bar two-cocycles. For the
comparison homotopies, enumerate finite global
quotients with ramification allowed in a finite
set $T$ containing the ramification of $o_n$ and
the support of $y_n$; enumerate larger finite
sets $T$ as well. Use the tame resolution at the
finite residue support. Thus only the arithmetic
representative is required to factor through an
$S$-ramified quotient; the global comparison
homotopy retains its necessary ramification.
Test the cochain
identities in the finite multiplication tables.
For the prescribed cycle a matching representative
and homotopy occur, so this search terminates.
The inverse in (C.15.5) is thereby a cochain
comparison, not a choice of a splitting of the
short exact sequence (C.12.1).
At a finite coefficient level the residue support
is finite: this is precisely where the restricted
condition in (C.15.1) is used. As with C.14, this
is an exact construction and no claim of practical
running time is made.

The group $G_{L,S}$ has three-cohomological dimension
at most two. Therefore $H^3(G_{L,S},\mu_{3^n})=0$,
which proves existence in (C.15.2). Its right side
is the standard evaluation cup product for
$A_n\otimes A_n^D\to\mu_{3^n}$. Since
$db_v=\operatorname{res}_v z$ and $d\xi_n=0$,
\[
 d\bigl((\operatorname{res}_v\xi_n)\cup b_v
                         -\operatorname{res}_v a\bigr)=0.
\]
We can obtain $a$ by enumerating finite $S$-ramified
quotients refining those through which $\xi_n,z$
factor and solving the finite linear cochain
equation. The vanishing just proved guarantees
termination of this search as well.

In the compact-support cone (C.14.8) the product
of $\xi_n$ and $(z,(b_v))$ is represented by
\[
 \bigl(\xi_n\cup z,
              ((\operatorname{res}_v\xi_n)\cup b_v)_v\bigr).
\]
Subtracting the coboundary of $(a,0)$ gives
$(0,((\operatorname{res}_v\xi_n)\cup b_v
                            -\operatorname{res}_v a)_v)$.
Normalize the compact-support trace by the sum of
the local Brauer invariants on such a representative.
This is the trace compatibility in Milne II.3,
preceding Theorem 3.1 \cite{MilneADT}.
The resulting trace is exactly (C.15.3). Each local
invariant is that of a finite crossed-product
algebra, so it is determined by finite extension,
norm, and valuation data over $L_v$.
The construction proves directly that (C.15.3)
is the perfect pairing (C.14.9), applied to the
class of $\xi_n$. It consequently descends to
$C_n$ and is independent of the representatives
and of $a$. In particular, changing $a$ by a
global two-cocycle changes the sum by zero, by
the global Brauer invariant relation.

At the integral level, injectivity of each $D_v$
proves that (C.15.1) is exactly the vanishing of
the residue obstruction, with its unique integral
correction. Under that condition the class of
$(o,y)$ is the unique arithmetic class of C.12.
By (C.14.5), this class vanishes exactly when
all its finite components pair to zero with
$C_n$. Those components are exactly (C.15.3).
This proves both directions of (C.15.4).
The finite-$S$ first-cohomology inverse-limit
identification and the stabilizer calculation
of C.12 give the stated torsor and its precise
conjugacy quotient. The global Milnor term is
still the one computed in F.1--F.3; the additional
finite arithmetic evaluations introduce no
new first derived limit. Thus the criterion
includes infinitely ramified classes, all
torsion classes, every prime above three, and
the modified real-place convention.
\end{proof}
